% ============================================================
% Step3 教程 · Part II：工程实现详解
% 由 step3_tutorial.tex \input 引入（共用同一份导言区与主题色）
% ============================================================

\part{工程实现详解}

\begin{tipbox}{这一部分写什么}
Part I 讲了知识；这一部分把仓库里\textbf{真实做出来的 Step3 工程}逐模块
讲透：每个参数为什么是这个值、每个赋值为什么用阻塞/非阻塞、
三轮 Kp 整定是怎么在闭环仿真里收敛的。对应目录：
\code{03\_vision/step3\_closed\_loop/}——2 个 RTL（209 行）、2 个 TB、
3 轮整定记录。
\end{tipbox}

% ============================================================
\section{工程总览}

\subsection{文件清单}

\begin{center}
\begin{tabular}{llr l}
\toprule
\rowcolor{rowtint}\textbf{文件} & \textbf{类型} & \textbf{行数} & \textbf{职责} \\
\midrule
\code{gimbal\_ctrl.v} & RTL & 140 & 控制律：死区+比例+限速+限幅，双轴 \\
\code{pwm\_servo.v} & RTL & 69 & 双路 50Hz PWM，共周期，上电中位 \\
\code{tb/tb\_gimbal.v} & 测试 & 62 & 控制律 10 步进表逐字段比对 \\
\code{tb/tb\_pwm.v} & 测试 & 94 & PWM 波形逐拍连续监视器（v2） \\
\code{sim/gimbal\_stim.txt} & 向量 & 10 & 帧输入序列 \\
\code{sim/gimbal\_exp.txt} & 向量 & 10 & 金标准期望 \\
\bottomrule
\end{tabular}
\end{center}

\subsection{数据流与时钟域}

\begin{Verilog}
Step2 检测链 (74.25MHz 像素域)
   |  目标坐标 + frame_tick (每帧 1 拍)
   v
pulse_sync  (单 bit 脉冲跨域, Step1 模块)
   v
gimbal_ctrl (50MHz PWM 域)  err/step/duty 计算
   |  O_duty_x/y_us (13bit, us)
   v
pwm_servo   (50MHz)         双路 50Hz 波形 -> 舵机
\end{Verilog}

两个时钟域：像素域 74.25MHz 跑检测，PWM 域 50MHz 跑控制与波形。
跨域只传两样东西：\textbf{帧拍脉冲}（pulse\_sync，单 bit）和
\textbf{目标坐标}（cdc\_4phase，多 bit）。50MHz 是刻意选择的：
1us 恰好 50 拍，µs 到拍的换算纯乘法无余数。

\subsection{两条设计原则（本步新增）}

\begin{keybox}{Step3 的方法论}
\textbf{1. 无符号纪律}：误差、乘积、累加一律显式 \code{\$signed}；
task 内中间量 29bit。整步最阴的 bug 全部出在符号上。\\
\textbf{2. 参数即固化}：所有控制参数是 RTL parameter——综合进 bit 流，
上电即生效，满足"冷启动直接跟踪"的比赛硬约束。
\end{keybox}

% ============================================================
\section{模块详解：gimbal\_ctrl}

\subsection{task 复用：双轴一份逻辑}

水平、垂直两轴的控制律完全同构，只有参数不同。RTL 用一个 task
\code{calc\_axis} 承载全部算术，时序块里调两次：

\begin{Verilog}
task calc_axis(input signed [12:0] err, input [15:0] kp, ...);
    prod = err * $signed({1'b0, kp});   // 29bit 有符号积
    s    = prod >>> 8;                  // Q8.8 -> us, floor
    ...
endtask
\end{Verilog}

注意 \code{\$signed(\{1'b0, kp\})}：位拼接默认产生无符号数，
不包一层 \code{\$signed} 的话整个乘法退化成无符号，负误差直接回绕——
这正是闭环发散事故的模型侧根因（详见振荡复盘一节）。

\subsection{五步顺序与实现}

\begin{Verilog}
dz_hit = (err >= -dz) && (err <= dz);      // 1) 死区判定
if (dz_hit) s = 0;                          //    命中则步进清零
if (s >  SLEW_P)  s =  SLEW_P;              // 2) 限速（双向）
else if (s < -SLEW_P) s = -SLEW_P;
nx = $signed({1'b0, duty_in}) + s;          // 3) 累加
if (nx > SMAX_P) begin duty_out = SMAX_P; lim_hit = 1; end
else if (nx < SMIN_P) begin duty_out = SMIN_P; lim_hit = 1; end
\end{Verilog}

\subsection{输出同拍锁存"新值"（一个真 bug）}

初版输出寄存写成 \code{O\_duty\_x\_us <= duty\_x}——而 \code{duty\_x}
在同一拍才被赋成新值，非阻塞赋值读的是\textbf{旧值}，输出滞后一帧。
tb\_gimbal 的步进表第一帧就对不上。修复：

\begin{Verilog}
O_duty_x_us <= (I_track_en && I_target_valid) ? dnx_x : duty_x;
//                              ^^^^^ 本拍刚算出的新值
\end{Verilog}

\begin{warnbox}{非阻塞赋值的"读旧写新"}
非阻塞赋值 RHS 立即求值、LHS 帧末更新。同一个 always 块里
"先更新 \code{duty\_x}，再输出 \code{duty\_x}"，读到的是更新前的值。
要么像本工程一样输出 \code{dnx\_x}（新值），要么想清楚就是要旧值。
\end{warnbox}

\subsection{状态字与丢失策略}

\begin{Verilog}
O_status <= {2'b00, limy, limx, lost, dzy, dzx};   // 一拍拼装
// lost 帧: duty 保持不动, O_lost_cnt 饱和累加
// track_en=0: duty 强制回中位, 状态清零
\end{Verilog}

一个迭代 1 抓出的细节：无效帧分支里 \code{dzx/limx} 清零用非阻塞，
同拍组装的 \code{O\_status} 采样到\textbf{旧值}——状态字残留上一帧的死区位。
修复：该分支的清零改\textbf{阻塞赋值}（ 组装在同一个顺序块里，
先清后拼）。这是"阻塞/非阻塞选择不是教条、而是看数据依赖"的实例。

% ============================================================
\section{模块详解：pwm\_servo}

\subsection{µs 到拍的换算}

\begin{Verilog}
localparam TICKS_PER_US = 50;               // 50MHz / 1e6
wire [19:0] pulse_x = duty_x * TICKS_PER_US; // 1500us -> 75000 拍
\end{Verilog}

20bit 计数器数到 1,000,000（50Hz 周期）回零，两路输出
\code{pwm = (period\_cnt < pulse)}——\textbf{一个计数器、两个比较器}，
共周期是结构保证的。

\subsection{第二道限幅与上电中位}

\begin{Verilog}
wire [12:0] dx = (I_duty_x_us > 13'd2500) ? 13'd2500 :
                 (I_duty_x_us < 13'd500)  ? 13'd500  : I_duty_x_us;
// 复位值: duty_x <= 13'd1500    —— 上电即中位
\end{Verilog}

控制律已经限在 700～2300，这里再限 500～2500——两道保险的区间
故意不同：如果上层逻辑跑飞给出荒谬值，第二道还能兜住机械极限。

% ============================================================
\section{测试平台设计}

\subsection{tb\_gimbal：10 步进表覆盖全部分支}

金标准 \code{gen\_tb\_gimbal()} 构造 10 帧输入序列，
\textbf{刻意覆盖每种截断}：

\begin{center}
\begin{tabular}{lll}
\toprule
\rowcolor{rowtint}\textbf{帧} & \textbf{输入} & \textbf{考察点} \\
\midrule
0-1 & cx=840 (err=200) & 常规比例步进（1500 到 1515） \\
2   & cx=900 (err=260) & 累加继续 \\
3   & cx=700 (err=60)  & 小误差小步进 \\
4   & cx=641 (err=1)   & \textbf{死区命中，duty 不动} \\
5-6 & valid=0          & \textbf{丢失保持 + lost\_cnt 累加} \\
7   & cx=100 (err=-540) & 负方向 + 大步进 \\
8   & cx=640 (err=0)   & 死区 \\
9   & cx=660 (err=20)  & 死区 \\
\bottomrule
\end{tabular}
\end{center}

期望文件每行 \code{\{dx:03X\}\{dy:03X\}\{lost:04X\}\{dzx:X\}\{dzy:X\}}，
TB 逐字段比对——duty、lost 计数、两个死区标志\textbf{都要对}。
第一条期望 \code{5EB5DC000001} 手工可验：dx=0x5EB=1515（err=200 步进 15）、
dy=0x5DC=1500（cy=360 正中死区）、dz\_y=1。

\subsection{tb\_pwm：从任务式同步到连续监视器}

第一版 TB 用任务式 \code{@(posedge tick)} 等待周期起点再测量——
实测 \code{len=1}：TB 的等待事件和 DUT 的置位\textbf{同拍竞态}，
采样点漂移。v2 重写为\textbf{连续监视器}：

\begin{Verilog}
// 自由运行计数器 + tick 锁存，没有任何 @ 事件同步
always @(posedge clk) begin
    if (pwm) pulse_len <= pulse_len + 1;
    if (tick) begin
        report(frame_cnt, pulse_len, period_len);
        pulse_len <= 0;
    end
    period_len <= (tick) ? 1 : period_len + 1;
end
\end{Verilog}

锁存窗口\textbf{包含 tick 拍本身}（v1.1 曾差一拍测得 999999），
换挡激励避开过渡周期。v2 测得：上电中位 75000/75000/1000000 逐拍精确，
换挡后 60000（1200us）/100000（2000us）精确。

\subsection{tb\_closed\_loop：相机-舵机模型在环}

真正的验收是把整条链闭环跑起来。TB 内建行为级相机模型
（参数都有实物依据）：

\begin{center}
\begin{tabular}{lll}
\toprule
\rowcolor{rowtint}\textbf{参数} & \textbf{取值} & \textbf{依据} \\
\midrule
光学放大 & 20 px/度 & 62.2 度 FOV @ 1280px \\
舵机增益 & 0.18 度/us & 180 度 / 1000us 行程 \\
机械速率限幅 & 1.5 度/帧 & 约 90 度/s \\
系统延迟 & 1 帧 & 检测链流水；控制节拍 = 新目标包到达拍（T1 修正前为 2 帧） \\
初始指向 theta0 & 30 度 & 与目标差 10 度 \\
\bottomrule
\end{tabular}
\end{center}

模型每个 PWM 周期算 \code{theta\_cmd = 30 + (duty-1500)×0.18}，
带速率限幅逼近；TPG 方块按 \code{x = 640 + (40-theta)×20} 实时移动
——\textbf{相机真的在"看"云台转完的画面}。验收标准：
60 帧内误差进入 ±25px 死区并保持 10 帧。

% ============================================================
\section{验证结果与振荡事故复盘}

\subsection{回归矩阵}

\begin{center}
\begin{tabular}{lll}
\toprule
\rowcolor{rowtint}\textbf{测试} & \textbf{比对粒度} & \textbf{结果} \\
\midrule
tb\_gimbal & 10 步进表逐字段（duty/lost/dz 标志） & PASS \\
tb\_pwm & 波形逐拍（脉宽/周期） & PASS \\
tb\_closed\_loop & 在环收敛（settle\_frame） & PASS：12 帧锁定 \\
\bottomrule
\end{tabular}
\end{center}

\subsection{Step3 抓到的真缺陷}

\begin{center}
\begin{tabular}{llp{5.4cm}l}
\toprule
\rowcolor{rowtint}\textbf{编号} & \textbf{现象} & \textbf{根因} & \textbf{类型} \\
\midrule
G1 & 输出滞后一帧 & 输出寄存同拍读旧值（非阻塞语义） & RTL \\
G2 & 状态字残留 & 无效帧清零用非阻塞，同拍拼装读旧值 & RTL \\
G3 & 测得脉宽 len=1 & 任务式 @ 同步与 DUT 置位同拍竞态 & TB \\
G4 & 周期 999999 & 计数窗口不含 tick 拍 & TB \\
G5 & 闭环跑飞 109 度 & 模型无符号减法回绕 + 回路增益超标 & 模型+参数 \\
\bottomrule
\end{tabular}
\end{center}

\subsection{振荡事故复盘（为什么值得单独一节）}

整定第一轮 Kp=90/70，闭环仿真 60 帧里角度\textbf{单调跑飞到 109 度}，
duty 打满限幅。复盘出两个叠加的问题：

\begin{warnbox}{事故 1：TB 模型的无符号减法}
相机模型算误差时用了无符号减法：\code{(40 - theta)} 在 theta 超过 40
时回绕成巨大正数，模型认为目标突然飞到另一边——闭环被喂了假误差。
修复：全链路 \code{\$signed}。教训：\textbf{行为模型也要遵守无符号纪律}，
它错的时候 RTL 怎么对都没用。\\[4pt]
\textbf{事故 2：回路增益超标}。Kp=90 时单帧增益 L=90×0.0144=1.30，
两帧延迟放大后 2L=2.60——每往返误差翻 2.6 倍，单调发散。
第二轮 Kp=40（2L=1.16）仍是慢发散；第三轮 Kp=20/18（2L=0.58）进入稳定收敛。
\textbf{闭环专项修正（T1）}：控制节拍改为"新目标包到达拍"，回路延迟由 2 帧降到
1 帧（判据放宽为 L<1），实测 \textbf{settle\_frame 由 36 改善到 12}，
RTL 默认参数同步更新。
\end{warnbox}

\begin{keybox}{给现场整定者的操作清单（完整版见 IF-SERVO-PWM.md）}
1）先验极性：给固定小误差，云台必须\textbf{朝减小误差}的方向动；
反了改 \code{INVERT\_X/Y} 参数，不要倒接电机线。\\
2）从保守 Kp 起步（按增益公式算出上限，取一半）。\\
3）逐步加大直到出现轻微振铃，回退 30\%。\\
4）SLEW 从 100 起步：只影响"最多多快"，不影响稳定性判据。\\
5）改完参数重新综合固化——参数是 RTL parameter，不烧不生效。
\end{keybox}

% ============================================================
\section{上板挂载说明（给上板组）}

\begin{center}
\begin{tabular}{ll}
\toprule
\rowcolor{rowtint}\textbf{项目} & \textbf{内容} \\
\midrule
时钟 & 50MHz PWM 域（官方例程已有该域，直接挂） \\
跨域输入 & 目标坐标（cdc\_4phase）；控制节拍本地生成 = 新目标包到达拍
（T1 后不再用 pulse\_sync 传 tick） \\
输出 & O\_pwm\_x/y 接舵机信号线；O\_duty\_dbg 供 ILA 观测 \\
按键/拨码 & track\_en 接拨码（0 = 回中位，演示/安全用） \\
供电 & 舵机独立 5V 供电，\textbf{与 FPGA 共地不共电} \\
参数 & KP/DZ/SLEW/限幅全部是 parameter，现场整定后重编译固化 \\
\bottomrule
\end{tabular}
\end{center}

上板三步：
\begin{enumerate}
  \item 两个 RTL 加入 TD 工程，实例化 gimbal\_ctrl + pwm\_servo；
  \item 把 Step2 的 \code{O\_frame\_tick}/坐标经跨域模块接入；
  \item 先不接舵机，用 ILA 看 O\_pwm 波形确认周期/脉宽，再带载。
\end{enumerate}

\begin{tipbox}{和 Step2 的分工}
Step2 只负责"看"，输出坐标和帧拍；Step3 只负责"动"，输入坐标和帧拍。
两者通过两个跨域信号解耦——检测链改坏了不影响舵机安全（丢失保持），
舵机故障也不阻塞检测。\textbf{接口即边界}，约定见 \code{06\_interfaces}。
\end{tipbox}

% ============================================================
\section{小结：Step3 的可复用产出}

\begin{itemize}
  \item \textbf{2 个可综合 RTL 模块}（209 行），控制律与波形生成解耦；
  \item \textbf{一套增益设计方法}：L = KP×0.0144，按回路延迟选判据
    （初版 2 帧延迟用 2L<1；T1 修正后延迟 1 帧，判据放宽为 L<1），
    换相机/舵机后重算即可；
  \item \textbf{一个在环验证范式}：行为级相机模型 + TPG 方块跟随，
    无硬件就能整定参数；
  \item \textbf{三份整定记录}：发散/慢发散/锁定的完整过程，
    是答辩时"工程严谨性"的最好素材。
\end{itemize}

\begin{okbox}{下一步（Step4 预告）}
坐标有了、云台动起来了，但屏幕上还看不到任何"系统在干活"的痕迹。
Step4 把结果画回画面：状态文字、目标框、跨域目标包、整链联调——
评委会看到的画面，从这里来。
\end{okbox}
